\documentclass[12pt,a4paper]{article}

% ================= PACKAGES =================
\usepackage[margin=1in]{geometry}
\usepackage[T1]{fontenc}
\usepackage{graphicx}
\usepackage{titlesec}
\usepackage{tabularx}
\usepackage{longtable}
\usepackage{array}
\usepackage{hyperref}
\usepackage{listings}
\usepackage{xcolor}
\usepackage{enumitem}
\usepackage{setspace}
\usepackage{mathptmx}

% ================= FONT SIZE DEFINITIONS =================
\newcommand{\sizeEighteen}{\fontsize{18pt}{22pt}\selectfont}
\newcommand{\sizeSixteen}{\fontsize{16pt}{20pt}\selectfont}
\newcommand{\sizeFourteen}{\fontsize{14pt}{17pt}\selectfont}
\newcommand{\sizeTwelve}{\fontsize{12pt}{15pt}\selectfont}

% ================= HEADER COLORS =================
\definecolor{RamaiahBlue}{RGB}{27,31,92}
\definecolor{RamaiahMaroon}{RGB}{139,30,63}
\definecolor{CodeBg}{RGB}{248,249,250}
\definecolor{CodeBorder}{RGB}{218,224,233}

% ================= LISTINGS FORMAT =================
\lstdefinestyle{codestyle}{
    backgroundcolor=\color{CodeBg},
    commentstyle=\color{green!50!black},
    keywordstyle=\color{blue!80!black}\bfseries,
    stringstyle=\color{orange!80!black},
    basicstyle=\ttfamily\fontsize{9.5pt}{11.5pt}\selectfont,
    numberstyle=\ttfamily\fontsize{8.5pt}{10.5pt}\selectfont\color{gray},
    breaklines=true,
    breakatwhitespace=false,
    frame=single,
    rulecolor=\color{CodeBorder},
    numbers=left,
    numbersep=6pt,
    xleftmargin=2.2em,
    aboveskip=4pt,
    belowskip=4pt,
    showstringspaces=false
}
\lstset{style=codestyle}

% ================= HEADING FORMATS =================
\titleformat{\section}{\sizeEighteen\bfseries}{\thesection}{1em}{}
\titleformat{\subsection}{\sizeSixteen\bfseries}{\thesubsection}{1em}{}
\titleformat{\subsubsection}{\sizeFourteen\bfseries}{\thesubsubsection}{1em}{}

\titlespacing*{\section}{0pt}{10pt}{6pt}
\titlespacing*{\subsection}{0pt}{8pt}{4pt}
\titlespacing*{\subsubsection}{0pt}{6pt}{3pt}

% ================= REPEATABLE HEADER BANNER =================
\newcommand{\reportheader}{%
\noindent
\begin{minipage}[c]{0.38\textwidth}
  \includegraphics[height=1.2cm]{ramaiah_logo.jpg}
\end{minipage}%
\begin{minipage}[c]{0.62\textwidth}
  \raggedleft
  {\color{RamaiahBlue}\fontsize{13pt}{16pt}\selectfont\bfseries Department of CSE (Artificial Intelligence \& Machine Learning) and CSE (Cyber Security)}
\end{minipage}\\[0.2cm]
\noindent\textcolor{RamaiahMaroon}{\rule{\linewidth}{1.2pt}}
\vspace{0.3cm}
}

% ================= INLINE FIGURE MACROS =================
\newcommand{\inlinefig}[3]{%
\noindent
\begin{minipage}{\linewidth}
\centering
\vspace{2pt}
\includegraphics[width=#2]{#1}\\[3pt]
{\fontsize{9.5pt}{11.5pt}\selectfont\textbf{#3}}
\vspace{4pt}
\end{minipage}
}
\newcommand{\inlinefigh}[4]{%
\noindent
\begin{minipage}{\linewidth}
\centering
\vspace{1pt}
\includegraphics[width=#2,height=#3,keepaspectratio]{#1}\\[2pt]
{\fontsize{9.5pt}{11.5pt}\selectfont\textbf{#4}}
\vspace{3pt}
\end{minipage}
}

\begin{document}
\pagestyle{empty}
\sizeTwelve

% ======================================================
% PAGE 1: TITLE PAGE
% ======================================================
\reportheader

\vspace{0.8cm}
\begin{center}
{\sizeEighteen \bfseries SELF LEARNING ACTIVITY REPORT \par}

\vspace{1.0cm}

{\sizeTwelve Submitted in partial fulfilment of the Requirements for \\[0.2cm]
the Academics of 5\textsuperscript{th} Semester, B.E in CSE (Artificial Intelligence and Machine Learning) \par}

\vspace{1.0cm}

{\sizeEighteen \bfseries DEVOPS \par}

\vspace{0.4cm}

{\sizeSixteen \bfseries 24CIAEC59 \par}

\vspace{1.0cm}

{\sizeTwelve In \par}

\vspace{0.3cm}

{\sizeFourteen \bfseries Dept. of CSE (Artificial Intelligence and Machine Learning) \par}

\vspace{1.0cm}

{\sizeTwelve By \par}

\vspace{0.3cm}

{\sizeSixteen \bfseries Gautam Rajesh \par}
\vspace{0.15cm}
{\sizeFourteen \bfseries USN: 1MS24CI045 \par}

\vspace{1.2cm}

{\sizeTwelve Under the guidance of \par}
\vspace{0.25cm}

{\sizeFourteen \bfseries Mr. Prashanth Reddy C \par}
\vspace{0.1cm}
{\sizeTwelve Senior Lecturer \par}
{\sizeTwelve Department of CSE (AI \& ML) and CSE (Cyber Security) \par}

\vspace{1.2cm}

{\sizeSixteen \bfseries M. S. Ramaiah Institute of Technology \par}
\vspace{0.25cm}
{\sizeTwelve (Autonomous Institute Affiliated to VTU) \\ MSR Nagar, MSRIT Post \par}
\vspace{0.1cm}
{\sizeTwelve Bangalore-560054, Karnataka, India \par}
\end{center}

\newpage

% ======================================================
% PAGE 2: CERTIFICATE
% ======================================================
\reportheader

\vspace{1.0cm}
\begin{center}
{\sizeEighteen \bfseries Certificate \par}
\end{center}

\vspace{1.0cm}

\noindent
{\sizeTwelve
This is to certify that the work carried out is a bonafide record of the project based learning activity work on the DevOps course done by \textbf{Gautam Rajesh (1MS24CI045)} under my supervision and guidance, in partial fulfilment of the requirements for the award of the Bachelor of Engineering in Computer Science \& Engineering (Artificial Intelligence and Machine Learning) from Ramaiah Institute of Technology, Bangalore for the academic year 2026-2027.
\par}

\vspace{4.0cm}

\noindent
{\sizeTwelve Candidate Sign \hfill Faculty Sign \par}

\newpage
\pagestyle{plain}

% ======================================================
% PAGE 3: LIST OF ACTIVITIES
% ======================================================
\section*{LIST OF PROJECT BASED LEARNING COMPONENTS}

\noindent
{\sizeTwelve
\textbf{COURSE NAME:} DevOps \hfill \textbf{COURSE CODE:} 24CIAEC59 \\
\textbf{CREDITS:} 1:0:0 \hfill \textbf{CONTACT HOURS:} 15L \\
\textbf{MAX. MARKS:} 20
\par}

\vspace{0.6cm}

{\sizeTwelve
\setlength{\tabcolsep}{4pt}
\renewcommand{\arraystretch}{1.25}
\begin{longtable}{|>{\centering\arraybackslash}m{1.2cm}|>{\raggedright\arraybackslash}p{10.6cm}|>{\centering\arraybackslash}m{1.8cm}|}
\hline
\multicolumn{1}{|c|}{\textbf{Sl. No}} &
\multicolumn{1}{c|}{\textbf{Description}} &
\multicolumn{1}{c|}{\textbf{Marks}} \\
\hline
\endhead

1 &
\textbf{CI/CD Pipeline using Git, Maven/Gradle and Jenkins.}
Set up a version-controlled build pipeline for a sample application. Create a Git repository, use feature branching and merging, configure a Maven or Gradle build with dependency management, integrate the project into a Jenkins job that pulls from Source Code Management (SCM), triggers automatically on commit, executes the build, performs a post-build action, and add one CI-related enhancement (automated tests, build versioning, or GitHub Actions).

\textit{Covers: Unit 1 (DevOps culture/pipeline), Unit 2 (Git, Maven/Gradle), Unit 3 (Jenkins CI). Maps to CO1, CO2, CO3.}
& 10 \\
\hline

2 &
\textbf{Configuration Management with Ansible and Deployment on Kubernetes.}
Write an Ansible playbook (using variables, handlers, and at least one role) to configure and deploy the application built in Component 1. Containerize the application and write Kubernetes manifests (Deployment, Service, and a ConfigMap/PVC where applicable). Deploy on a local cluster (Minikube/Kind), demonstrate scaling using replica sets, show the use of labels and selectors, and add one deployment-related enhancement (rollback, ConfigMap, probes, multi-stage Dockerfile, second microservice, or health endpoint).

\textit{Covers: Unit 4 (Ansible), Unit 5 (Kubernetes). Maps to CO4, CO5.}
& 10 \\
\hline

\end{longtable}
\par}

\newpage

% ======================================================
% PAGE 4: ASSESSMENT RUBRICS
% ======================================================
\section*{Assessment Rubrics}
{\sizeTwelve Evaluation for \textbf{each} 10-mark component will be based on the following rubrics.\par}

\vspace{0.4cm}

{\sizeTwelve
\setlength{\tabcolsep}{4pt}
\renewcommand{\arraystretch}{1.25}
\begin{longtable}{|>{\raggedright\arraybackslash}p{3.0cm}|>{\raggedright\arraybackslash}p{3.5cm}|>{\raggedright\arraybackslash}p{3.5cm}|>{\raggedright\arraybackslash}p{3.5cm}|}
\hline

\textbf{Criteria} & \textbf{2 --- Good} & \textbf{1 --- Average} & \textbf{0 --- Poor} \\

\hline
\endhead
Concept Understanding (2 marks) & Answers technical design and architecture questions correctly and immediately regarding CI triggers, build lifecycles, and orchestration concepts & Demonstrates basic understanding but requires prompting on deeper lifecycle mechanisms & Lacks conceptual understanding or cannot explain design choices \\
\hline
Core Implementation (2 marks) & Automated pipeline executes completely end-to-end without manual intervention or errors & Pipeline functions but requires minor workarounds or manual triggering & Pipeline fails to execute or only static artifacts are presented \\
\hline
Enhancement (2 marks) & Innovation and enhancement clearly demonstrated live and integrated cleanly into pipeline & Enhancement attempted but partially functional or not fully integrated & No enhancement implemented \\
\hline
Analysis \& Interpretation (2 marks) & Accurately analyzes build logs, runtime metrics, and scaling dynamics with clear rationale & Describes operational outcomes without analyzing underlying system behavior & Does not interpret output or explain results \\
\hline
Documentation \& Presentation (2 marks) & Comprehensive, personalized documentation with clear architecture diagrams and exact terminal logs & Complete report with minor formatting gaps or omitted outputs & Incomplete documentation or mismatched evidence \\
\hline
\end{longtable}
\par}

\newpage

% ======================================================
% PAGE 5: MARKS AWARDED
% ======================================================
\section*{Marks Awarded}

\vspace{0.3cm}

{\sizeTwelve
\setlength{\tabcolsep}{3pt}
\renewcommand{\arraystretch}{1.5}
\begin{longtable}{|>{\centering\arraybackslash}p{1.0cm}|>{\raggedright\arraybackslash}p{2.2cm}|>{\raggedright\arraybackslash}p{2.2cm}|>{\raggedright\arraybackslash}p{2.0cm}|>{\raggedright\arraybackslash}p{2.3cm}|>{\raggedright\arraybackslash}p{2.5cm}|>{\centering\arraybackslash}p{1.4cm}|}
\hline
\textbf{Sl. No} &
\textbf{Concept Understanding (2)} &
\textbf{Core Implementation (2)} &
\textbf{Enhancement (2)} &
\textbf{Analysis \& Interpretation (2)} &
\textbf{Documentation \& Presentation (2)} &
\textbf{Total} \\
\hline
\parbox[c][1.1cm][c]{1.0cm}{\centering 1} & & & & & & \\
\hline
\parbox[c][1.1cm][c]{1.0cm}{\centering 2} & & & & & & \\
\hline
\multicolumn{6}{|c|}{\parbox[c][1.1cm][c]{13.0cm}{\centering\textbf{TOTAL MARKS}}} &
\parbox[c][1.1cm][c]{1.4cm}{\centering} \\
\hline
\end{longtable}
\par}

\newpage

% ======================================================
% PAGE 6: COMPONENT 1 (START)
% ======================================================
\section*{Component 1: CI/CD Pipeline using Git, Maven/Gradle and Jenkins}

\subsection*{Abstract}
This component builds an automated continuous integration pipeline for \textbf{CloudCart}, an enterprise e-commerce inventory and order processing microservice built with Java 17 and Spring Boot 3.3.5. Source code is maintained in Git using structured feature branches merged via non-fast-forward commits. Apache Maven manages compilation, dependencies, and fat JAR packaging. A Jenkins declarative pipeline polls GitHub for changes, runs five automated JUnit 5 tests, archives the packaged JAR, and tags a Docker image with the build number and commit hash.

\subsection*{Introduction}
Continuous Integration accelerates software delivery by integrating and verifying developer changes frequently. Git provides distributed version control so multiple developers work safely on independent branches. Maven standardizes project construction, handles dependency resolution, and packages runnable artifacts. Jenkins unifies these tools into an automated feedback loop that pulls code upon each push, executes tests, and reports code quality before release.

In this activity, the application is \textbf{CloudCart}, exposing REST endpoints: \texttt{GET /api/products}, \texttt{POST /api/orders}, \texttt{GET /api/info}, and \texttt{GET /health}. The service manages real inventory limits: purchasing an item when stock is exhausted returns HTTP 409 Conflict, providing a robust business logic condition for automated unit test validation.

\subsection*{Implementation}

\subsubsection*{Development Environment}
\begin{itemize}[nosep]
    \item OS: Ubuntu 22.04 LTS on WSL2 / Linux VM
    \item Git 2.34.1, remote repository on GitHub (\texttt{Gautam-Rajesh/cloudcart})
    \item OpenJDK 17.0.12 for application compilation; Jenkins 2.452 LTS
    \item Apache Maven 3.9.6 with Spring Boot 3.3.5 and JUnit 5 / MockMvc
    \item Jenkins 2.452 LTS with Git, Pipeline, and SCM polling plugins
    \item Sample app: CloudCart E-Commerce Order Microservice
\end{itemize}

\subsubsection*{Step 1: Git Repository and Branching}
Development was completed on two feature branches merged into \texttt{main} using \texttt{--no-ff} to record clear merge nodes in history.

\begin{lstlisting}[language=bash]
git init cloudcart && cd cloudcart
git add . && git commit -m "Initial CloudCart microservice skeleton"

git checkout -b feature/info-endpoint
# implement /api/info (application version, config message, pod name)
\end{lstlisting}

\newpage

% ======================================================
% PAGE 7: STEP 1 CONT & STEP 2 MAVEN POM.XML
% ======================================================
\begin{lstlisting}[language=bash,firstnumber=6]
git commit -am "Add /api/info endpoint backed by config and pod name"
git checkout main
git merge --no-ff feature/info-endpoint -m "Merge feature/info-endpoint"

git checkout -b feature/devops
# add Jenkinsfile, GitHub Actions, Dockerfile, k8s manifests, Ansible role
git checkout main
git merge --no-ff feature/devops -m "Merge feature/devops"

git remote add origin https://github.com/Gautam-Rajesh/cloudcart.git
git push -u origin main
\end{lstlisting}

\subsubsection*{Step 2: Maven Build Configuration (pom.xml excerpt)}
The project inherits from \texttt{spring-boot-starter-parent} 3.3.5 to manage dependency versions centrally. Setting \texttt{finalName} to \texttt{cloudcart} produces a predictable target artifact \texttt{target/cloudcart.jar}; the \texttt{spring-boot-maven-plugin} packages it as an executable fat JAR.

\begin{lstlisting}[language=XML]
<parent>
    <groupId>org.springframework.boot</groupId>
    <artifactId>spring-boot-starter-parent</artifactId>
    <version>3.3.5</version>
</parent>
<groupId>com.rit</groupId>
<artifactId>cloudcart</artifactId>
<version>1.0.0</version>
<properties>
    <java.version>17</java.version>
</properties>
<dependencies>
    <dependency>
        <groupId>org.springframework.boot</groupId>
        <artifactId>spring-boot-starter-web</artifactId>
    </dependency>
    <dependency>
        <groupId>org.springframework.boot</groupId>
        <artifactId>spring-boot-starter-validation</artifactId>
    </dependency>
    <dependency>
        <groupId>org.springframework.boot</groupId>
        <artifactId>spring-boot-starter-test</artifactId>
        <scope>test</scope>
    </dependency>
</dependencies>
<build>
    <finalName>cloudcart</finalName>
    <plugins>
\end{lstlisting}

\newpage

% ======================================================
% PAGE 8: STEP 2 CONT & STEP 3 JENKINSFILE
% ======================================================
\begin{lstlisting}[language=XML,firstnumber=24]
        <plugin>
            <groupId>org.springframework.boot</groupId>
            <artifactId>spring-boot-maven-plugin</artifactId>
        </plugin>
    </plugins>
</build>
\end{lstlisting}

\subsubsection*{Step 3: Jenkins Pipeline (Jenkinsfile)}
The declarative pipeline lives in source control. The Jenkins job is configured as \textit{Pipeline script from SCM} pointing to \texttt{https://github.com/Gautam-Rajesh/cloudcart.git}, branch \texttt{*/main}. The \texttt{pollSCM} trigger polls GitHub every two minutes, starting builds upon detecting newly pushed commits.

\begin{lstlisting}[language=bash]
pipeline {
    agent any
    tools { maven 'Maven3'; jdk 'JDK17' }
    options { disableConcurrentBuilds() }
    triggers { pollSCM('H/2 * * * *') }
    environment { IMAGE = "cloudcart" }

    stages {
        stage('Checkout') { steps { checkout scm } }
        stage('Build') { steps { sh 'mvn -B clean compile' } }
        stage('Test') {
            steps { sh 'mvn -B test' }
            post { always { junit 'target/surefire-reports/*.xml' } }
        }
        stage('Package') { steps { sh 'mvn -B package -DskipTests' } }
        stage('Archive') {
            steps { archiveArtifacts artifacts: 'target/*.jar', fingerprint: true }
        }
        stage('Docker Build') {
            steps {
                script {
                    env.GIT_SHA = sh(script: 'git rev-parse --short HEAD', returnStdout: true).trim()
                }
                sh 'docker build -t ${IMAGE}:${BUILD_NUMBER}-${GIT_SHA} .'
            }
        }
    }
\end{lstlisting}

\newpage

% ======================================================
% PAGE 9: STEP 3 POST & STEP 4 ENHANCEMENT & DESCRIPTION
% ======================================================
\begin{lstlisting}[language=bash,firstnumber=27]
    post {
        success { echo "Built ${IMAGE}:${BUILD_NUMBER}-${GIT_SHA}" }
        failure { echo 'Build failed' }
    }
}
\end{lstlisting}

\subsubsection*{Step 4: CI Enhancement}
Two CI enhancements were implemented and demonstrated:
\begin{itemize}[nosep]
    \item \textbf{Automated tests in CI}: A dedicated Test stage executes five JUnit 5 / MockMvc tests (\texttt{healthReturnsUp}, \texttt{listsProducts}, \texttt{orderReturns201ThenOutOfStockReturns409}, \texttt{unknownProductReturns404}, \texttt{orderRejectsInvalidQuantity}) prior to packaging, publishing results via the \texttt{junit} step to render test trend metrics in Jenkins.
    \item \textbf{Build versioning}: The Docker image is tagged \texttt{\$\{BUILD\_NUMBER\}-\$\{GIT\_SHA\}} instead of static tags, ensuring absolute traceability back to the Jenkins build and Git commit.
    \item A GitHub Actions workflow (\texttt{.github/workflows/ci.yml}) running \texttt{mvn -B clean package} on every push provides redundant cloud verification.
\end{itemize}

\subsubsection*{Description of the Activity}
The Jenkins job synchronizes directly with GitHub and executes stages sequentially: checkout, compile, unit test, package, archive, and Docker build. Build \#1 took roughly 40 seconds to download dependencies and base layers; subsequent builds completed in 28 seconds due to caching. Build \#1 was initiated manually to confirm configuration, and Build \#3 was triggered from the dashboard. Build \#4 was triggered without manual interaction: after pushing the commit \textit{"Live push demo: update banner message"}, the SCM polling log reported \textit{Changes found} and Jenkins printed \textit{Started by an SCM change}, proving automated commit triggering.

The post-build action archived \texttt{cloudcart.jar} (19.42 MB) on the job page. The test trend graph showed five passing tests across every build, and uniquely tagged images were generated (\texttt{cloudcart:1-7b3e01a}, \texttt{2-4f81c92}, \texttt{3-e29a3b1}, \texttt{4-de82f19}).

During initial configuration, Jenkins daemon service issues occurred when running on incompatible JDK packages; this was resolved by configuring JDK 17 explicitly in Jenkins Global Tool Configuration while keeping pipeline options managed directly in the repository \texttt{Jenkinsfile}.

\newpage

% ======================================================
% PAGE 10: RESULTS COMPONENT 1 (FIG 1 - 4)
% ======================================================
\subsection*{Results}

\inlinefigh{fig1_git_history.png}{0.92\linewidth}{2.8cm}{Figure 1: Git history: main, feature/info-endpoint and feature/devops, with explicit merge commits c4d92e1 and 7b3e01a}

\inlinefigh{fig2_app_startup.png}{0.92\linewidth}{2.7cm}{Figure 2: CloudCart starting as an executable JAR (java -jar target/cloudcart.jar) on port 9091}

\inlinefigh{fig3_order_curl.png}{0.92\linewidth}{1.8cm}{Figure 3: First order returns 201 Created; second order for last stock unit returns 409 Conflict}

\inlinefigh{fig4_maven_test.png}{0.92\linewidth}{1.8cm}{Figure 4: mvn clean package: Tests run: 5, Failures: 0, BUILD SUCCESS}

\newpage

% ======================================================
% PAGE 11: RESULTS COMPONENT 1 (FIG 5 - 7)
% ======================================================
\inlinefig{fig5_jenkins_console.png}{0.92\linewidth}{Figure 5: Jenkins build \#1 console: image tagged cloudcart:1-7b3e01a, Finished: SUCCESS}

\vspace{0.3cm}

\inlinefig{fig6_jenkins_stages.png}{0.92\linewidth}{Figure 6: Jenkins Stages view: every stage of build \#1 is green}

\vspace{0.3cm}

\inlinefig{fig7_jenkins_build3.png}{0.92\linewidth}{Figure 7: Build \#3: archived cloudcart.jar, git revision, tests with no failures, and the commits included}

\newpage

% ======================================================
% PAGE 12: RESULTS COMPONENT 1 (FIG 8 - 9)
% ======================================================
\inlinefig{fig8_scm_polling.png}{0.92\linewidth}{Figure 8: SCM polling log reporting Changes found after a new commit reached GitHub}

\vspace{0.8cm}

\inlinefig{fig9_jenkins_build4.png}{0.92\linewidth}{Figure 9: Build \#4 console: Started by an SCM change, the automatic trigger}

\newpage

% ======================================================
% PAGE 13: RESULTS COMPONENT 1 (FIG 10 & CONCLUSION)
% ======================================================
\inlinefig{fig10_jenkins_dashboard.png}{0.92\linewidth}{Figure 10: Job page after four builds: archived artifact, test trend graph (5 tests passed in every build) and build history}

\vspace{0.6cm}

\subsection*{Conclusion}
Building this pipeline demonstrated that the core benefit of Continuous Integration is rapid, automated verification: once the push, poll, test, archive, and container tagging workflow was operational, code adjustments automatically generated verified artifacts without manual Jenkins interactions. Positioning automated testing before the packaging phase guaranteed that business rule regressions stop flawed binaries from entering production.

\newpage

% ======================================================
% PAGE 14: COMPONENT 2 (START)
% ======================================================
\section*{Component 2: Configuration Management with Ansible and Deployment on Kubernetes}

\subsection*{Abstract}
This component takes the CloudCart JAR produced in Component 1 and deploys it across two environments. An Ansible role installs Java, sets up a dedicated service user, copies the executable JAR, and manages it as a systemd service with dynamic port overrides and notification handlers. The application is containerized using a multi-stage Dockerfile and deployed to a local Minikube cluster as a 3-replica Deployment behind a NodePort Service, featuring externalized ConfigMap parameters, health probes, horizontal scaling, and rollback validation.

\subsection*{Introduction}
Once built, software artifacts must be reliably deployed and orchestrated. Ansible provides agentless configuration management using declarative YAML playbooks to describe desired machine states, ensuring idempotent host setup. Kubernetes orchestrates containers across clusters, handling automated scheduling, replica scaling, service discovery, and self-healing.

In this activity, Ansible configures the local Ubuntu host using local connection, and Minikube with the Docker driver provides the Kubernetes v1.37.0 cluster runtime.

\subsection*{Implementation}

\subsubsection*{Development Environment}
\begin{itemize}[nosep]
    \item Ansible (apt package) on Ubuntu 22.04 / WSL2, managed host: localhost
    \item Docker Engine, multi-stage build using \texttt{maven:3.9-eclipse-temurin-17} and \texttt{eclipse-temurin:17-jre}
    \item Minikube v1.39.0 with Docker driver (2 CPUs, 1900 MB, bridge CNI), Kubernetes v1.37.0, containerd 2.3.4
    \item kubectl 1.35.10 CLI
\end{itemize}

\subsubsection*{Step 1: Ansible Inventory, Playbook and Role}
The role was scaffolded using \texttt{ansible-galaxy init roles/cloudcart}. The inventory specifies the \texttt{appservers} group, and the playbook applies the role with privilege escalation.

\begin{lstlisting}[language=bash]
# inventory
[appservers]
localhost ansible_connection=local

# playbook.yml
---
- name: Set up CloudCart runtime
  hosts: appservers
  become: true
  roles:
    - cloudcart
\end{lstlisting}

\newpage

% ======================================================
% PAGE 15: ANSIBLE DEFAULTS & TASKS
% ======================================================
Role variable defaults (\texttt{roles/cloudcart/defaults/main.yml}); \texttt{app\_port} is overridden at runtime with \texttt{-e app\_port=9090}:

\begin{lstlisting}[language=bash]
---
app_port: 8080
app_dir: /opt/cloudcart
app_user: cloudcart
jar_src: "{{ playbook_dir }}/../target/cloudcart.jar"
app_message: "CloudCart deployed by Ansible"
\end{lstlisting}

Role tasks (\texttt{roles/cloudcart/tasks/main.yml}):

\begin{lstlisting}[language=bash]
---
- name: Install Java 17 runtime
  ansible.builtin.apt:
    name: openjdk-17-jre-headless
    state: present
    update_cache: true

- name: Create application user
  ansible.builtin.user:
    name: "{{ app_user }}"
    system: true
    shell: /usr/sbin/nologin

- name: Create application directory
  ansible.builtin.file:
    path: "{{ app_dir }}"
    state: directory
    owner: "{{ app_user }}"
    mode: "0755"

- name: Copy application jar
  ansible.builtin.copy:
    src: "{{ jar_src }}"
    dest: "{{ app_dir }}/cloudcart.jar"
    owner: "{{ app_user }}"
    mode: "0644"
  notify: Restart cloudcart

- name: Install systemd unit
  ansible.builtin.template:
    src: cloudcart.service.j2
    dest: /etc/systemd/system/cloudcart.service
    mode: "0644"
  notify:
    - Reload systemd
    - Restart cloudcart
\end{lstlisting}

\newpage

% ======================================================
% PAGE 16: ANSIBLE HANDLERS, TEMPLATE, & DOCKERFILE
% ======================================================
\begin{lstlisting}[language=bash,firstnumber=38]
- name: Enable and start service
  ansible.builtin.systemd:
    name: cloudcart
    enabled: true
    state: started
    daemon_reload: true
\end{lstlisting}

Handlers (\texttt{roles/cloudcart/handlers/main.yml}) and systemd template (\texttt{templates/cloudcart.service.j2}):

\begin{lstlisting}[language=bash]
---
- name: Reload systemd
  ansible.builtin.systemd:
    daemon_reload: true

- name: Restart cloudcart
  ansible.builtin.systemd:
    name: cloudcart
    state: restarted

# cloudcart.service.j2
[Unit]
Description=CloudCart E-Commerce App
After=network.target

[Service]
User={{ app_user }}
Environment=SERVER_PORT={{ app_port }}
Environment=APP_MESSAGE={{ app_message }}
ExecStart=/usr/bin/java -jar {{ app_dir }}/cloudcart.jar
Restart=on-failure

[Install]
WantedBy=multi-user.target
\end{lstlisting}

Run command:
\begin{lstlisting}[language=bash]
cd ansible
ansible-playbook -i inventory playbook.yml -K -e app_port=9090
\end{lstlisting}

\subsubsection*{Step 2: Dockerfile (multi-stage)}
The first stage compiles the code with Maven and JDK; the second stage copies only the packaged JAR into a slim JRE image, ensuring build dependencies do not inflate runtime image size.

\begin{lstlisting}[language=bash]
# ---- build stage ----
FROM maven:3.9-eclipse-temurin-17 AS build
WORKDIR /app
COPY pom.xml .
\end{lstlisting}

\newpage

% ======================================================
% PAGE 17: DOCKERFILE CONT & K8S MANIFESTS
% ======================================================
\begin{lstlisting}[language=bash,firstnumber=5]
RUN mvn -B dependency:go-offline
COPY src ./src
RUN mvn -B clean package -DskipTests

# ---- runtime stage ----
FROM eclipse-temurin:17-jre
WORKDIR /app
COPY --from=build /app/target/cloudcart.jar app.jar
EXPOSE 8080
ENTRYPOINT ["java", "-jar", "app.jar"]
\end{lstlisting}

\subsubsection*{Step 3: Kubernetes ConfigMap, Deployment and Service}
The Deployment runs three replicas labelled \texttt{app: cloudcart}. The Service selector matches the Pod template labels so traffic routes across all replicas. Configuration parameters are injected via ConfigMap using \texttt{envFrom}, and probes query \texttt{/health}.

\begin{lstlisting}[language=bash]
# k8s/configmap.yaml
apiVersion: v1
kind: ConfigMap
metadata:
  name: cloudcart-config
data:
  SERVER_PORT: "8080"
  APP_MESSAGE: "Welcome to CloudCart (from ConfigMap)"
  APP_VERSION: "1.0.0"

# k8s/deployment.yaml
apiVersion: apps/v1
kind: Deployment
metadata:
  name: cloudcart
spec:
  replicas: 3
  selector:
    matchLabels:
      app: cloudcart
\end{lstlisting}

\newpage

% ======================================================
% PAGE 18: K8S CONT, STEP 4 & STEP 5
% ======================================================
\begin{lstlisting}[language=bash,firstnumber=21]
  template:
    metadata:
      labels:
        app: cloudcart
    spec:
      containers:
        - name: cloudcart
          image: cloudcart:1.0
          ports:
            - containerPort: 8080
          envFrom:
            - configMapRef:
                name: cloudcart-config
          readinessProbe:
            httpGet: { path: /health, port: 8080 }
            initialDelaySeconds: 10
            periodSeconds: 5
          livenessProbe:
            httpGet: { path: /health, port: 8080 }
            initialDelaySeconds: 20
            periodSeconds: 10

# k8s/service.yaml
apiVersion: v1
kind: Service
metadata:
  name: cloudcart-svc
spec:
  type: NodePort
  selector:
    app: cloudcart
  ports:
    - port: 80
      targetPort: 8080
      nodePort: 30080
\end{lstlisting}

\subsubsection*{Step 4: Deploy and Scale}
\begin{lstlisting}[language=bash]
docker build -t cloudcart:1.0 .
minikube start --driver=docker --memory=1900mb --cpus=2 --cni=bridge
minikube image load cloudcart:1.0
kubectl apply -f k8s/
kubectl get deployments,services,pods -o wide
minikube service cloudcart-svc --url

kubectl scale deployment cloudcart --replicas=5
kubectl get pods
kubectl scale deployment cloudcart --replicas=3
kubectl get pods
\end{lstlisting}

\subsubsection*{Step 5: Deployment Enhancement}
Several deployment enhancements were implemented; the first two were demonstrated live:
\begin{itemize}[nosep]
    \item \textbf{Rollback}: A deliberately broken image was deployed with \texttt{kubectl set image deployment/cloudcart cloudcart=cloudcart:broken}, and \texttt{kubectl rollout undo deployment/cloudcart}
\end{itemize}

\newpage

% ======================================================
% PAGE 19: STEP 5 CONT, DESCRIPTION & FIG 11 (TREE)
% ======================================================
restored the working version.
\begin{itemize}[nosep]
    \item \textbf{ConfigMap}: The port, welcome message, and version are read from \texttt{cloudcart-config} through environment variables instead of hard-coded values; \texttt{GET /api/info} returns them together with the active Pod name.
    \item \textbf{Probes, multi-stage Dockerfile and health endpoint}: Readiness and liveness probes target \texttt{/health}, which returns status and uptime in seconds; the Dockerfile is multi-stage as shown above.
\end{itemize}

\subsubsection*{Description of the Activity}
The Ansible role is idempotent. On the run shown below, the JAR copy task reported \texttt{changed}, notifying the \textit{Restart cloudcart} handler, while tasks already in the desired state (Java installation, user, directory, systemd unit) reported \texttt{ok}. Handlers execute only when modified state warrants a daemon restart. Passing \texttt{-e app\_port=9090} overrode the default role port, and querying \texttt{/health} on port 9090 returned \texttt{UP}, proving that runtime variables applied cleanly via the systemd template.

On Kubernetes, \texttt{kubectl apply -f k8s/} created the ConfigMap, Deployment, and Service. Three Pods progressed from \texttt{ContainerCreating} to \texttt{Running} and transitioned to \texttt{1/1} ready once the readiness probe succeeded. The Service exposes NodePort 30080, distributing traffic across replicas: successive calls to \texttt{/api/info} returned distinct Pod hostnames. Scaling to five spawned two new replicas, and scaling back down cleanly terminated extra instances while existing Pods maintained uptime. During the rollback test, the broken image entered \texttt{ErrImagePull}, yet existing replicas continued handling live traffic without outage. Executing \texttt{rollout undo} restored the Deployment to the working ReplicaSet immediately.

Network delays when downloading default CNI images were avoided by starting Minikube with \texttt{--cni=bridge}, and memory allocation was constrained to 1900 MB to operate efficiently on the host system.

\subsection*{Results}

\inlinefig{fig11_ansible_tree.png}{0.88\linewidth}{Figure 11: Role structure generated by ansible-galaxy init (defaults, handlers, tasks, templates)}

\newpage

% ======================================================
% PAGE 20: RESULTS COMPONENT 2 (FIG 12 - 13)
% ======================================================
\inlinefig{fig12_ansible_run.png}{0.92\linewidth}{Figure 12: Ansible run with -e app\_port=9090: failed=0, handler Restart cloudcart executed, and /health and /api/products answering on port 9090}

\vspace{0.4cm}

\inlinefig{fig13_docker_build.png}{0.92\linewidth}{Figure 13: Multi-stage Docker build, images tagged by Jenkins (cloudcart:1-7b3e01a to 4-de82f19), and container responding on /health}

\newpage

% ======================================================
% PAGE 21: RESULTS COMPONENT 2 (FIG 14 - 15)
% ======================================================
\inlinefig{fig14_k8s_startup.png}{0.92\linewidth}{Figure 14: minikube cluster started and application manifests applied}

\vspace{0.4cm}

\inlinefig{fig15_k8s_resources.png}{0.92\linewidth}{Figure 15: Pods created, then kubectl get deployments,services,pods -o wide: 3/3 ready, selector app=cloudcart, NodePort service 80:30080}

\newpage

% ======================================================
% PAGE 22: RESULTS COMPONENT 2 (FIG 16 - 17 & CONCLUSION)
% ======================================================
\inlinefig{fig16_k8s_scaling.png}{0.92\linewidth}{Figure 16: Scaling from 3 to 5 replicas and back to 3, with /api/info answered by different Pods}

\vspace{0.2cm}

\inlinefig{fig17_k8s_rollback.png}{0.92\linewidth}{Figure 17: Rollback: a broken image gives ErrImagePull, kubectl rollout undo recovers, and the ConfigMap message is returned by /api/info}

\vspace{0.2cm}

\subsection*{Conclusion}
This component highlighted the architectural differences between bare-metal configuration and container orchestration. Ansible made server provisioning repeatable and deterministic through idempotent tasks and conditional handlers. Kubernetes demonstrated declarative self-healing and horizontal scalability. The rollback experiment showed how rolling updates isolate failures and preserve application availability, while ConfigMaps decouple environment configuration from immutable container images.

\newpage

% ======================================================
% PAGE 23: REFERENCES
% ======================================================
\section*{References}

\begin{enumerate}[nosep]
    \item Git Version Control System Documentation -- \url{https://git-scm.com/doc}
    \item Apache Maven Build Lifecycle and Plugin Guides -- \url{https://maven.apache.org/guides/}
    \item Jenkins Automation Server User Handbook -- \url{https://www.jenkins.io/doc/}
    \item Ansible Community Automation Documentation -- \url{https://docs.ansible.com/ansible/latest/}
    \item Kubernetes Production Container Orchestration -- \url{https://kubernetes.io/docs/}
    \item Spring Boot Reference Documentation -- \url{https://spring.io/projects/spring-boot}
    \item Docker Engine and Multi-Stage Container Packaging -- \url{https://docs.docker.com/}
\end{enumerate}

\end{document}
